\documentclass[10pt,a4paper]{article}
\usepackage[fleqn]{amsmath}
\usepackage{amssymb,amsthm,graphicx}
\usepackage{latexsym}
\usepackage[fleqn]{amsmath}
\usepackage{amssymb}
\usepackage{amsthm}
\usepackage{enumitem}
\setlist{itemsep=0.4pt,topsep=0pt}
\usepackage{hyperref}
\newcommand{\HCF}{\mathrm{hcf}}
\newcommand{\ZZ}{\mathbb{Z}}
\newcommand{\rk}{\mathrm{rk}}
\DeclareMathOperator{\len}{length}
\DeclareMathOperator{\bl}{b}
\DeclareMathOperator{\cl}{cl}
\DeclareMathOperator{\di}{d}
\DeclareMathOperator{\Cl}{Cl}
\DeclareMathOperator{\acl}{acl}
\DeclareMathOperator{\tp}{tp}
\DeclareMathOperator{\Th}{Th}
\newcommand{\Cor}{\mathcal{C}_\mathrm{or}}
\newcommand{\Cun}{\mathcal{C}_\mathrm{un}}
\newcommand{\mons}{\mathbb{M}}
\newcommand{\strong}{\mathrm{strong}}
\newcommand{\axioms}{\Delta_{\mathrm{FSTS}}}

%\definecolor{magenta}{gray}{0}

\newtheoremstyle{note}% hnamei
{8pt}% hSpace abovei
{14pt}% hSpace belowi
{\normalfont}% hBody fonti
{}% hIndent amounti1
{\bfseries}% hTheorem head fonti
{ }% hPunctuation after theorem headi
{.5em}% hSpace after theorem headi2
{}% hTheorem head spec (can be left empty, meaning ‘normal’)i
\theoremstyle{note}
\newtheorem{thm}{Theorem}[section]
%\newtheorem{theorem}{Theorem}
\newtheorem{lemma}[thm]{Lemma} 
\newtheorem{df}[thm]{Definition}
\newtheorem{rmk}[thm]{Remark} 
\newtheorem{prop}[thm]{Proposition}
\newtheorem{fact}[thm]{Fact}
\newtheorem{cor}[thm]{Corollary}
\newtheorem{conj}[thm]{Conjecture}
\newtheorem{ex}[thm]{Example}
\newtheorem{qu}[thm]{Question}
\newtheorem*{notation}{Notation} 
\title{Quantitative n-amalgamation for finite unconfined Steiner triple systems}
\author{Silvia Barbina, Enrique Casanovas}
\date{arxiv:2411.13723v3; CC BY 4.0}
\makeatletter
\newlabel{unconfined}{{source: unconfined}{}{}{}{}}
\makeatother
\begin{document}
\noindent\textit{Reference locators: references to unnumbered external source sections retain their original author labels; complete source files accompany this item.}\par
\maketitle
\noindent\textit{Editorial note.} C\_un is the class of finite unconfined partial Steiner triple systems (source437-440). The whole counted task has finite structures; n-strongness is the bounded test-set condition686-692.\par\medskip
\noindent\textit{Editorial note.} The original author construction and proof are retained. The relation A<=B is strong-substructure order, not the HF linear order itself.\par\medskip
\noindent\textit{Editorial note.} The external support proposition is by the same two authors in Endomorphisms of free Steiner quasigroups, fixed version2412.06055v3. Its finite induction supplies the criterion needed here; the printed infinite compactness clause includes a noted a!=b omission, but the counted finite construction does not use that clause.\par\medskip
\noindent\textit{Editorial note.} In the support source line293, the greatest point of a finite subset A-prime is said to belong to at most one block of A. The needed and directly justified assertion is at most one block inside A-prime: the other two points of each such block lie below that greatest point in the HF-order. The printed broader intermediate wording is retained; the subset restriction follows from the source definition and is the conclusion used here.\par\medskip
\noindent\textit{Editorial scope.} One quantitative n-amalgamation proposition for finite unconfined partial Steiner systems. Its identification process and local strongness lemmas are solution obligations. The second source provides the actual HF-ordering criterion. Infinite-model homogeneity, stability, generic existence and the previously published tree lemma are outside this task.\par\medskip
\noindent\textit{Allowed standard prerequisites (editorial).}\begin{itemize}
\item Finite partial structures and embeddings. An embedding is an isomorphism onto the induced partial structure on its image.. Track the two copies and their intersection during point identification. No amalgamation theorem is assumed.
\item Finite induction and propositional compactness. For the included generic HF-order criterion; the actual application here is finite and uses the finite vertex-removal induction.
\end{itemize}\medskip
A  partial Steiner triple system can be regarded as a first-order structure in the language containing a ternary relation symbol $R$ and satisfying the axioms
\begin{itemize}
\item for all permutations $\sigma \in S_3$, $ \forall x_1 x_2 x_3 \, \left(R(x_1, x_2, x_3) \rightarrow R(x_{\sigma(1)}, x_{\sigma(2)}, x_{\sigma(3)} \right)$
\item $\forall x y z \left( R(x,y,z) \rightarrow y\neq z \right)$
\item $\forall x y w z \left( R(x,y,w) \wedge R(x,y,z) \rightarrow w=z \right)$.
\end{itemize}
The relation $R$ is interpreted as the set of blocks: three elements $x, y$ and $z$ are in a block if and only if $R(x,y,z)$.  A partial STS that satisfies the additional axiom
\begin{itemize}
\item $\forall x y \left( x \neq y \rightarrow \exists z \,R(x,y,z)\right)$
\end{itemize}
is an STS.
With this choice of language, a substructure of an STS is a partial STS.

An STS can also be described as a \textbf{Steiner quasigroup}, that is, a structure with a single binary operation symbol $\cdot$, interpreted as an operation that satisfies the following axioms:
 \begin{enumerate}
\item $\forall x y \, (x \cdot y= y\cdot x)$
\item $\forall x \, (x\cdot x =x)$
\item $\forall x y \, (x\cdot (x\cdot y)=y)$.
\end{enumerate}
Thus, Steiner quasigroups are an equational class --- see \cite{ganter1, ganter2}. This is what makes it possible to employ universal algebraic tools.

A quasigroup $M$ determines an STS whose blocks are the elements of the graph of the operation $\cdot$ without the diagonal, that is,
\[ \mathcal{B} = \{ \{x,y,z \} \subseteq M : x \cdot y = z \} \smallsetminus \{ \{x \} : x \in M \} \, .\]
It is easy to see that each STS determines a quasigroup and each quasigroup determines an STS, where
  $$x \cdot y = z \mbox{ iff } R(x,y,z) \mbox{ or } x=y=z.$$  In fact, as  noted in \cite{ganter1, ganter2}, this shows that each Steiner triple system is bi-interpretable with a Steiner quasigroup, and justifies our choice to use the term \textit{Steiner triple system} to cover both the functional and relational structures from the end of this section onwards.
  %shown in \cite{baldwin23}, Section~3, each Steiner triple system is bi-interpretable with a Steiner quasigroups.  

	\begin{df} A \textbf{hyperfree ordering}  (or \textbf{HF-ordering}) of a partial Steiner triple system  $A$ is a linear ordering $<$ of $A$ such that for every $a\in A$ there is at most one pair of smaller elements $\{a_1,a_2\}\subseteq A$ with $a=a_1\cdot a_2$,  that is, there is at most one block $\{a,a_1,a_2\}$ with $a_1,a_2<a$.
	\end{df}

				\begin{df} %A partial STS $A\neq \emptyset$ is \textbf{unconfined}, or \textbf{contains no confined configuration}, if for all finite $A^\prime \subseteq A$ there is $a \in A^\prime$ that belongs to at most one block in $A^\prime$, that is, $a$ is {\color{magenta} either not a product of any elements in $A^\prime$, or the product of exactly two elements of $A^\prime$.}
 A finite partial STS  $A$ is \textbf{confined} if every element of $A$ belongs to at least two blocks.

A partial STS %$A\neq \emptyset$ 
 $A$ is \textbf{unconfined}, or \textbf{contains no confined configuration}, if no finite $A^\prime \subseteq A$ is confined.

Equivalently,  $A$ is unconfined if for all finite  nonempty $A^\prime \subseteq A$ there is $a \in A^\prime$ that belongs to at most one block in $A^\prime$, that is, $a$ is either not a product of any elements in $A^\prime$, or the product of exactly one pair of elements of $A^\prime$.
\end{df}
Note that being unconfined is stronger than not being confined: an unconfined partial STS is hereditarily not confined. 

	\begin{df} \label{def5_1} For $A, B \in \Cun$,
$$A \leq B \  \Leftrightarrow \ A\subseteq B \mbox{ and there is an HF-ordering of $B$ where $A$ is an initial segment}. $$
When $A \leq B$, we say that $A$ is \textbf{strong} in $B$.
\end{df}

Since $A$ and $B$ are finite, the condition $A \leq B$ is first-order.

To see this, we define certain formulas in the relational language of Steiner triple systems (so blocks interpret a ternary relation $R$ that is symmetric and irreflexive). In what follows, we implicitly assume that the elements $a_i$ are pairwise distinct. There is a formula $\delta_1(x_1, \ldots, x_n)$, in the relational language of partial STSs, that says that $x_1<x_2<\ldots <x_n$ is an HF-ordering, namely
$$\delta_1(x_1, \ldots, x_n) \equiv \bigwedge_{i=1}^n \bigg[ \bigwedge_{1\leq j <k<i} \bigg( R(x_i, x_j, x_k) \rightarrow \bigwedge_{\tiny \begin{array}{l} 1\leq j' <k'<i \\ \{j,k\} \neq \{j',k'\} \end{array}} \neg R(x_i,x_{j'},x_{k'}) \bigg) \bigg] \, .$$
%Then $\delta_1(a_1, \ldots, a_n)$ holds if and only if $a_1<a_2<\ldots <a_n$ is an HF-ordering.

Now let $S_n$ denote the symmetric group on $n$ elements, and let
$$\delta_2 (x_1, \ldots, x_n) \equiv \bigvee_{\pi \in S_n} \delta_1(x_{\pi(1)}, \ldots, x_{\pi(n)}).$$
Then $\delta_2(a_1, \ldots, a_n)$ holds if and only if there is an HF-ordering of $a_1, \ldots, a_n$.

Consider the following modification of $\delta_2$:
$$\delta_3 (x_1, \ldots, x_n, x_{n+1}, \ldots , x_{n+m}) \equiv \bigvee_{\footnotesize \begin{array}{l} \pi \in S_n \\ \rho \in S_m \end{array}} \delta_1(x_{\pi(1)}, \ldots, x_{\pi(n)}, x_{n+\rho(1)}, \ldots, x_{n+\rho(m)}).$$
Then for $A=\{a_1, \ldots, a_n\} \subseteq \{a_1, \ldots, a_n, a_{n+1}, \ldots, a_{n+m}\}=B$ we have 
\[ B \models \delta_3 (a_1, \ldots, a_n, a_{n+1}, \ldots , a_{n+m}) \mbox{ if and only if } A\leq B \,. \]
So for each $m,n$, $A \subseteq B$ with  $|A|=n$ and $|B|=n+m$, there is a formula that asserts that $A \leq B$.

\begin{df}  Let $A,B\in \Cun$ and $n<\omega$.   We write  $A\leq^n B$  if  $A \subseteq B$  and for every  $X\subseteq B\smallsetminus A$  such that  $|X|\leq n$, we have  $A\leq A\cup X$. 

Given  a partial unconfined STS $B$, possibly infinite, and $A\in \Cun$ such that   $A\subseteq B$, we write  $A\leq^n B$  if $A\leq^n B_0$ for every  finite $B_0\subseteq B$ . 

In either case we say that $A$ is \textbf{$n$-strong} in $B$. 
\end{df}
Note that $A\leq B$ if and only if $A\leq^n B$  for every $n$. 

\begin{df}
	We say that $(\Cun, \leq)$  has the \textbf{$n$-amalgamation property} if  for every $A\leq B$ in $\Cun$ there is some $m<\omega$  such that whenever $C\in \Cun$ and $A\leq^m C$, there are $D\in \Cun$ and embeddings $f:B\rightarrow D$ and $g:C\rightarrow D$   such that 
$$ f\restriction A= g\restriction A, \quad f(B)\leq^n D \, \mbox{ and } \, g(C)\leq D \, .$$
\end{df}

	\begin{prop}\label{Feb_23_2024_2} $(\Cun,\leq)$ has the $n$-amalgamation property for every $n$.  In fact, if $A\leq B$ and $n<\omega$ are given, we can always take $m= n+|B\smallsetminus A|$.
	\end{prop}

\begin{prop}\label{2024_nov3} The following are equivalent for any partial Steiner triple system $A$:
\begin{enumerate}
\item $A$ is unconfined.
\item There is an HF-ordering of $A$.
\end{enumerate}	
\end{prop}	
\begin{proof} Consider first the case $A$ finite. By induction on $|A|$ it is easy to prove that, if $A$ is unconfined, then it has an HF-ordering. On the other hand, if there is an HF-ordering  of $A$, for any non-empty $A' \subseteq A$ the greatest element $a\in A'$ belongs to  at most one block of $A$, hence to at most a block of $A'$. Consider now the general case. If $A$ has an HF-ordering, every finite subset of $A$ has an HF-ordering and so it is unconfined, which implies that $A$ is unconfined. For the other direction, we use propositional compactness. Assume that every finite subset  of $A$ has an HF-ordering. For every pair  $(a,b)\in A\times A$, let $p_{ab}$ be a corresponding propositional variable. The set of sentences  
$$ \{p_{ab}\wedge p_{bc}\rightarrow p_{ac}\mid a,b,c\in A\}\cup \{\neg p_{aa}\mid a\in A\}\cup \{p_{ab}\vee p_{ba}\mid a,b\in A\}$$
 defines a linear ordering on $A$. If we add 
 $$\{\neg (p_{ba}\wedge p_{ca}\wedge p_{b^\prime a}\wedge p_{c^\prime a} )\mid \{a,b,c\}\neq\{a,b^\prime,c^\prime\} \text{ blocks of } A\}$$ 
 we are defining an HF-ordering on $A$. Every finite subset is satisfiable and compactness gives a valuation that satisfies the whole set.
\end{proof}

\begin{lemma}\label{diamond} Let $A,A_1,A_2 \in \Cun$. Then
	\begin{enumerate}
		\item $\emptyset \leq A$  
	\item If $B\in \Cun$ is such that $A_1\subseteq B$ and $A_2\leq B$, then $A_1\cap A_2\leq A_1$.
	\end{enumerate}
\end{lemma}

\begin{proof} 1. By Proposition~\ref{unconfined},  since every $A\in \Cun$ has an HF-ordering.
	
	2. Given an HF-ordering of $B$ with initial segment $A_2$, the elements of $A_1$ inherit an HF-ordering with initial segment $A_1\cap A_2$.
\end{proof}

\begin{lemma}\label{sat3}  If  $A\leq^n B$ and $B\leq^n C$, then $A\leq^n C$.
\end{lemma}
\begin{proof}  Assume $X\subseteq C \smallsetminus A$  and $|X|\leq n$. We check that  $A\leq A\cup X$. Since $B\leq^n C$, we have $B\leq B\cup X$. By Lemma~\ref{diamond},   $A\cup (X\cap B)\leq A\cup X$.
	Since $A\leq^n B$, we know that $A\leq A\cup (X\cap B)$.  Hence  $A\leq A\cup X$.
\end{proof}
	

	
\begin{lemma}\label{Feb_23_2024_1}\begin{enumerate}
			\item Assume $A\leq B$, $ b\in B\smallsetminus A$ and there is a block $\{a_1,a_2,b\}$  with  $a_1,a_2\in A$. Then  $Ab\leq B$.
			\item Assume $A\leq^n B$, $b\in B\smallsetminus A$ and there is a block $\{a_1,a_2,b\}$  with  $a_1,a_2\in A$. Then  $Ab\leq^{n-1} B$.
		\end{enumerate}
	\end{lemma}
	\begin{proof} 1.  Start with an HF-ordering of $B$ with initial segment $A$. Then $\{a_1,a_2,b\}$ is the only block containing $b$ and elements  smaller than $b$ in the order.  We can make $b$ the least element in $B\smallsetminus A$ and obtain, again, an HF-ordering of $B$, now with $Ab$ as initial segment. 
		
		2. Similar proof.
	\end{proof}

	\begin{proof} Assume $A\leq B$  and  $A\leq^m C$, where $m= n+|B\smallsetminus A|$. We want to find $B^\ast\cong_A B$ and $D\in \Cun$  such that  $C\leq D$ and $B^\ast\leq^n D$.  We may assume $C\cap B=A$. Note that the blocks containing only elements of $A$ are the same in $B$ and in $C$.  Assume there are $b\in B\smallsetminus A$, $c\in C\smallsetminus A$ and  $a_1,a_2\in A$  such that $\{a_1,a_2,b\}$ is a block of $B$ and  $\{a_1,a_2,c\}$ is a block of $C$. Redefine $A^\prime= A\cup\{c\}$,  $B^\prime = (B\smallsetminus \{b\})\cup \{c\}$ and modify the blocks of $B$ by exchanging $b$ with $c$  to obtain the blocks of $B^\prime$. Then $B^\prime\cong_A B$ and $B^\prime\cap C=A^\prime$.  By Lemma~\ref{Feb_23_2024_1}, we have  $A^\prime\leq B^\prime$ and $A^\prime\leq^{m-1} C$. Iterating this process gives $B^\ast\cong_A B$  such that if $A^\ast =B^\ast \cap C$, then   $A^\ast \leq B^\ast$,  $A^\ast\leq^{n +i}C$  for some $i \geq 0$  and there are no $b\in B^\ast\smallsetminus A^\ast$, $c\in C\smallsetminus A^\ast$ and $a_1,a_2\in A^\ast$ such that  $\{a_1,a_2,b\}$ is a block of $B^\ast$  and  $\{a_1,a_2,c\}$ is a block of $C$.  Now consider $D=B^\ast\cup C$ as a free amalgam of $B^\ast$ and $C$ over $A^\ast$, that is,  with the blocks in $B^\ast$ and in $C$ and no others.  Then $D$ is a partial STS. We claim that $D\in \Cun$, $C\leq D$ and $B^\ast\leq^n D$.
		
First we show that $D$ has an HF-ordering where $C$ is an initial segment. This proves that $D\in\Cun$ and $C\leq D$.  Take an HF-ordering $<_C$ of $C$ and an HF-ordering $<_{B^\ast}$ of  $B^\ast=A^\ast\cup (B^\ast\smallsetminus C)$ where $A^\ast$ is an initial segment. Then the sum 
$$(C,<_C)+(B^\ast\smallsetminus C,<_{B^\ast}\restriction (B^\ast\smallsetminus C))$$ 
is an HF-ordering of $D$. If not, there is an element of $B^\ast\smallsetminus C$ that belongs to different blocks containing smaller elements, but these elements then are in $A^\ast\cup (B^\ast\smallsetminus C)$, contradicting the choice of the ordering $<_{B^\ast}$ of $B^\ast$ as an HF-ordering over $A^\ast$.
		
		Finally, we check that $B^\ast\leq^n D$.  Let $X\subseteq D\smallsetminus B^\ast = C\smallsetminus A^\ast$  with $|X|\leq n$; we show that $B^\ast\leq B^\ast \cup X$.  Since  $A^\ast\leq B^\ast$, we have an HF-ordering $<_{B^\ast}$ of $B^\ast$ with initial segment $A^\ast$.  Since $A^\ast\leq^n C$, we know that $A^\ast\leq A^\ast\cup X$, and hence there is an HF-ordering $<_{A^\ast \cup X}$ of $A^\ast\cup  X$ with initial segment $A^\ast$. We combine these orderings by taking the sum $(B^\ast,<_{B^\ast})+(X,<_{A^\ast\cup X}\restriction X)$. Since an element of $X$ cannot have a block containing elements of $B^\ast$ unless the elements belong to $A^\ast$, we obtain an HF-ordering of $B^\ast\cup X$ with initial segment $B^\ast$.
	\end{proof}

\begin{thebibliography}{9}
\bibitem{barcas2}
S.~Barbina and E.~Casanovas.
\newblock Endomorphisms of free {S}teiner quasigroups, 2024.
\newblock{\url{arXiv:2412.06055}}.

\bibitem{casanovas}
E.~Casanovas.
\newblock Compactly expandable models and stability.
\newblock {\em The Journal of Symbolic Logic}, 60:673--683, 1995.
\bibitem{ganter1}
B.~Ganter and H.~Werner.
\newblock Equational classes of {S}teiner systems.
\newblock {\em Algebra Universalis}, 5:125--140, 1975.

\bibitem{ganter2}
B.~Ganter and H.~Werner.
\newblock Co-ordinatizing {S}teiner systems.
\newblock In C.C. Linder and A.~Rosa, editors, {\em Topics on {S}teiner
  systems}, pages 3--24. North Holland, 1980.
\end{thebibliography}
\end{document}
